\section{Validity and limits of generalization}
\label{sec:heldout-limits}
The primary study tests one requested model alias, one reasoning setting and one CLI version on a sampled function-generation benchmark. The repeated development stage does not supply model-level replication of the main sample. Public task availability leaves training contamination unresolved; related tasks can also violate an independent-task approximation. A different model, language, task distribution or service revision could reverse an observed contrast. Recorded dates and runtime hashes identify the executed client, not immutable provider weights.

The deadline and administrative continuation create a second boundary. Failed generations remain missing, and the amendment is disclosed. Inference from available pairs concerns the observed eligible subset. Assignment-denominator ranges show extreme missing-outcome possibilities without correcting selection bias. A timeout could reflect computational difficulty rather than random transport noise, so failure to reject a quality difference does not prove equivalent quality. No externally justified quality-loss margin or joint quality--cost criterion supports a non-inferiority or economic-superiority decision.

The interventions concern explicit instructions. Role labels do not certify independent agents, and call boundaries also expose intermediate responses. Although task information and stage operations are controlled, there is no matched hard output-token allowance. Resource comparisons include induced reasoning length, repeated context transmission and client overhead. List-price valuation depends on cache state; the no-cache sensitivity removes the discount but does not make subscription execution identical to an API deployment. Research and evaluator compute are excluded and no total cost of ownership is estimated.

Native tests improve auditability while retaining benchmark limitations. Gold and incorrect controls detect an unusable environment or an insensitive negative check, but passing those controls does not prove that every original assertion implements the natural-language task. A separate development audit records instruction/test disagreements as annotations, with original prompts, tests and scores unchanged. Those annotations are not post-hoc exclusions or a replacement gold standard. The seven main reference-ineligible tasks likewise remain in the assignment and resource accounting.

The INoT comparison is a disclosed prompt-level adaptation with an explicit resolution of source ambiguity. Its separate batch retains timing and cache differences. Self-Collaboration has two completed development checks on three previously exposed tasks, including nine fresh SCC/SR/SN assignments. Its role-excised variant and the planned 1,000-task external comparison remain unexecuted. Native author INoT, Agentless and SWE-agent implementations are also not evaluated. The latter systems introduce retrieval, tools and repository interaction outside this fixed-context factorial question. The one-issue SWE-bench demonstration provides genuine patch and test evidence, but its complete ten-assignment matrix on a single issue cannot support a repository-benchmark score or comparative agent ranking. A native multi-repository comparison and a matched-budget, multi-model replication remain necessary for those broader claims.
